% Ad Atticum 12.34 (ad-atticum-12-34)
% Source: https://raw.githubusercontent.com/PerseusDL/canonical-latinLit/master/data/phi0474/phi057/phi0474.phi057.perseus-lat1.xml
% Fetched by scripts/fetch_latin.py

% Dateline (Perseus): Scr. Asturae iii K. Apr. a. 709 (45).

\ciceroLetterOpener{CICERO ATTICO salutem}

\ciceroSection{1} ego hic vel sine Sicca (Tironi enim melius est) facillime possem esse ut in malis sed, quom scribas videndum mihi esse ne opprimar, ex quo intellegam te certum diem illius profectionis non habere, putavi esse commodius me istuc venire; quod idem video tibi placere. cras igitur in Siccae suburbano. inde , quem ad modum suades, puto me in Ficulensi fore.

\ciceroSection{2} quibus de rebus ad me scripsisti, quoniam ipse venio, coram videbimus. tuam quidem et in agendis nostris rebus et in consiliis ineundis mihique dandis in ipsis litteris quas mittis benevolentiam, diligentiam, prudentiam mirifice diligo. tu tamen si quid cum Silio, vel illo ipso die quo ad Siccam venturus ero, certiorem me velim facias, et maxime cuius loci detractionem fieri velit. quod enim scribis extremi, vide ne is ipse locus sit cuius causa de tota re, ut scis, est a nobis cogitatum. Hirti epistulam tibi misi et recentem et benevole scriptam.
